\documentclass[11pt,a4paper]{article}

\usepackage{amsmath}
\usepackage{mathtools}
\usepackage{fontspec}
\usepackage{unicode-math}
\setmathfont{KpMath-Regular.otf}
\setmathfont{KpMath-Bold.otf}[version=bold]
\usepackage{microtype}
\usepackage[dvipsnames]{xcolor}
\usepackage[margin=24mm, top=28mm, bottom=28mm]{geometry}

\usepackage{fancyhdr}
\pagestyle{fancy}
\fancyhf{}
\fancyhead[L]{\small\itshape The Plane-Wave Response of a Sphere}
\fancyhead[R]{\small\itshape 19 September 2026}
\fancyfoot[C]{\thepage}
\renewcommand{\headrulewidth}{0.4pt}

\setcounter{secnumdepth}{3}
\usepackage{booktabs}
\usepackage{array}
\usepackage[font=small, labelfont=bf, labelsep=period]{caption}
\usepackage{setspace}
\setstretch{1.12}
\setlength{\parskip}{0.4em}
\usepackage[colorlinks=true, linkcolor=blue!60!black, urlcolor=blue!60!black]{hyperref}

\newcommand{\bI}{\mathbf{I}}
\newcommand{\bR}{\mathbf{R}}
\newcommand{\bT}{\mathbf{T}}
\newcommand{\bu}{\mathbf{u}}
\newcommand{\bk}{\mathbf{k}}
\newcommand{\br}{\mathbf{r}}
\newcommand{\bta}{\boldsymbol{\tau}}
\newcommand{\bsig}{\boldsymbol{\sigma}}
\newcommand{\bep}{\mathbf{e}}
\newcommand{\khat}{\hat{\bk}}
\newcommand{\rhat}{\hat{\br}}
\newcommand{\zhat}{\hat{\mathbf{z}}}
\newcommand{\kperp}{\bk_\perp}
\newcommand{\tp}{^{\mkern-1.5mu\mathsf{T}}}

\title{\bfseries The Plane-Wave Response of a Sphere\\[4pt]
\large The angular spectrum of the elastic Mie field, and the reflection and\\
transmission matrices on planes bounding the scatterer}
\author{}
\date{19 September 2026}

\begin{document}
\maketitle

\begin{abstract}
\noindent
A homogeneous elastic sphere in a homogeneous whole space scatters an incident
plane wave into a field whose exact solution is a series of spherical waves.
This note gives that field's \emph{plane-wave} content: the amplitude, at each
lateral wavenumber, of every up- and down-going $P$, $SV$ and $SH$ wave crossing
a plane placed just outside the sphere.  Those amplitudes are the reflection and
transmission matrices $\bR(\kperp)$ and $\bT(\kperp)$.

The construction is closed-form.  Each outgoing multipole has a plane-wave
representation whose angular weight is its own spherical harmonic evaluated at
the plane-wave direction, so the spectrum of the scattered field is read off its
Mie coefficients with no spatial transform, no lateral window and no truncation.
Everything needed to implement it is given: the potentials and their
coefficients, the displacement and traction, the plane-wave representation and
its branch, the three channel spectra with their polarisations, the quadrature
path that removes the branch-point singularity, and the assembly into $\bR$ and
$\bT$.

Conventions are fixed in \S\ref{sec:conv} and used throughout: time dependence
$e^{-i\omega t}$, outgoing radial functions $h_n^{(1)}$, and $\operatorname{Im}
k_z \ge 0$.
\end{abstract}

\tableofcontents

%=====================================================================
\section{Statement of the problem}\label{sec:problem}

A sphere of radius $a$, with Lam\'e parameters $\lambda_1,\mu_1$ and density
$\rho_1$, sits in an unbounded homogeneous medium $\lambda,\mu,\rho$ with wave
speeds
\begin{equation}
  \alpha = \sqrt{(\lambda+2\mu)/\rho},\qquad \beta = \sqrt{\mu/\rho},
  \qquad k_P = \omega/\alpha, \qquad k_S = \omega/\beta .
\end{equation}
A plane wave is incident along $+z$.  Take two planes $z = -z_p$ and
$z = +z_p$ with $z_p > a$, so that both lie wholly outside the sphere.

We want, at each lateral wavenumber $\kperp = (k_x,k_y)$, the amplitudes of the
plane waves crossing those planes.  Above the sphere every scattered wave
travels in $-z$; below, in $+z$.  Collecting the three polarisations gives a
$3\times3$ reflection matrix on the upper plane and a $3\times3$ transmission
matrix on the lower one, each a function of $\kperp$.

\subsection{Conventions}\label{sec:conv}

\begin{itemize}
\item Time dependence $e^{-i\omega t}$; outgoing spherical waves carry
      $h_n^{(1)}(kr)$.
\item $z$ increases downward and is the first Cartesian index, so a vector is
      ordered $(z,x,y)$.  Spherical angles are measured from $+z$:
      $\cos\theta = z/r$, and $\varphi$ is the azimuth in the $(x,y)$ plane.
\item The vertical wavenumber of a plane wave is taken on the branch
      \begin{equation}\label{eq:branch}
        k_z(\kperp) = \sqrt{k^2 - |\kperp|^2}, \qquad \operatorname{Im}k_z \ge 0,
      \end{equation}
      real inside the light circle $|\kperp| < k$ and positive-imaginary outside
      it, so that $e^{i k_z |z|}$ is outgoing in both half-spaces: propagating
      within the circle and decaying beyond it.
\item Seismic units are used throughout --- $\mathrm{km/s}$,
      $\mathrm{g/cm^3}$, $\mathrm{GPa}$, $\mathrm{km}$.  In SI the matrices of
      \S\ref{sec:rt} carry displacement rows of order unity against traction
      rows of order $\mu k \sim 10^{6}$, and appear ill conditioned for a
      reason that is purely one of units.
\end{itemize}

%=====================================================================
\section{The exact field in spherical waves}\label{sec:mie}

\subsection{Potentials}

Write the displacement through the Helmholtz decomposition
\begin{equation}\label{eq:helm}
  \bu = \nabla\phi \;+\; \nabla\times\nabla\times(\br\psi) \;+\; \nabla\times(\br\chi),
\end{equation}
where $\br$ is the position vector and each potential satisfies a scalar
Helmholtz equation: $\phi$ with $k_P$, and $\psi$, $\chi$ with $k_S$.  The three
terms are the $L$-, $N$- and $M$-type families, radiating $P$, $SV$ and $SH$
respectively.  The two shear terms are divergence free, so
\begin{equation}\label{eq:div}
  \nabla\cdot\bu = \nabla^2\phi = -k_P^2\,\phi .
\end{equation}

Outside the sphere the scattered potentials are outgoing:
\begin{equation}\label{eq:pots}
  \phi = \sum_n a_n\,h_n^{(1)}(k_P r)\,P_n^m(\cos\theta)e^{im\varphi},
  \qquad
  \psi = \sum_n b_n\,h_n^{(1)}(k_S r)\,P_n^m(\cos\theta)e^{im\varphi},
\end{equation}
and likewise $\chi$ with coefficients $c_n$.  Inside, the same forms with
$j_n$ and the interior wavenumbers.

An axial $P$ wave is azimuthally symmetric, so it excites only $m = 0$, and it
excites no toroidal field: $c_n = 0$.  Incident $SV$ and $SH$ excite $m = 1$.
The channel formulas of \S\ref{sec:spectra} cover both, with the angular
functions \eqref{eq:angexplicit} giving the $m=0$ and $m=1$ cases explicitly.

\subsection{The incident field}\label{sec:incident}

The incident wave enters twice, in different forms, and both are needed.

\paragraph{As a series, to drive the boundary conditions.}  A unit-amplitude
$P$ wave travelling along $+z$ is $\bu^{\rm inc} = \zhat\,e^{ik_Pz}$, which
comes from the potential $\phi^{\rm inc} = e^{ik_Pz}/(ik_P)$ since
$\bu = \nabla\phi$.  The Rayleigh expansion
\begin{equation}\label{eq:rayleigh}
  e^{ikz} = \sum_{n\ge0}(2n+1)\,i^{\,n} j_n(kr)\,P_n(\cos\theta)
\end{equation}
then gives the incident partial-wave coefficients
\begin{equation}\label{eq:aninc}
  a^{\rm inc}_n = \frac{(2n+1)\,i^{\,n}}{i k_P},
  \qquad b^{\rm inc}_n = c^{\rm inc}_n = 0 ,
\end{equation}
so axial $P$ incidence is purely $m=0$ and purely $L$-type.

A unit shear wave along $+z$ is transverse, $\bu^{\rm inc} = \hat{\mathbf{x}}
e^{ik_Sz}$ for $SV$ and $\hat{\mathbf{y}}e^{ik_Sz}$ for $SH$.  Its expansion
carries $P_n^1$ and $e^{\pm i\varphi}$, so shear incidence is $m=1$, with the
$N$-type and $M$-type potentials respectively.  The conversion from the $m=0$
plane-wave coefficients to the $m=1$ ones is
\begin{equation}\label{eq:renorm}
  \varrho(n) = -\,\frac{1}{n(n+1)} ,
\end{equation}
which carries the normalisation of $P_n^1$ relative to $P_n$.

\paragraph{As a plane wave, in the angular spectrum.}  The same incident field
is, trivially, a single plane wave, so its spectrum is a delta:
\begin{equation}\label{eq:incspec}
  \hat{\bu}^{\rm inc}(\kperp) = (2\pi)^2\,\bep_\iota\,\delta^2(\kperp - \kperp^0),
\end{equation}
with $\kperp^0 = \mathbf{0}$ for axial incidence.  It travels in $+z$, so it is
present on the lower plane and absent on the upper one.  That is the whole of
the difference between $\bR$ and $\bT$ in \S\ref{sec:rt}: reflection is purely
scattered, transmission is the incident delta plus the scattered part.

\subsection{Coefficients}

The $a_n,b_n$ follow from continuity of $u_r$, $u_\theta$, $\sigma_{rr}$ and
$\sigma_{r\theta}$ at $r=a$ --- a $4\times4$ system per order, with columns
(scattered $P$, scattered $S$, interior $P$, interior $S$) and right-hand side
built from \eqref{eq:aninc}, or its shear counterpart through
\eqref{eq:renorm}.  The $c_n$ follow from the corresponding $2\times2$ system
in $u_\varphi$ and $\sigma_{r\varphi}$.  These are the classical Mie relations
and are not restated here.

\subsection{Displacement}

For $m=0$ the $L$- and $N$-type fields have no azimuthal component, and with
$z_n \equiv h^{(1)}_n$,
\begin{align}
  \text{$L$-type:}\quad
    & u_r = a_n\,k_P\,z_n'(k_P r)\,P_n, &
    & u_\theta = a_n\,\frac{z_n(k_P r)}{r}\,\frac{dP_n}{d\theta}, \label{eq:uL}\\[2pt]
  \text{$N$-type:}\quad
    & u_r = b_n\,n(n+1)\,\frac{z_n(k_S r)}{r}\,P_n, &
    & u_\theta = b_n\!\left[\frac{z_n(k_S r)}{r} + k_S z_n'(k_S r)\right]\!\frac{dP_n}{d\theta}. \label{eq:uN}
\end{align}
The $M$-type is purely azimuthal,
\begin{equation}\label{eq:uM}
  u_r = u_\theta = 0, \qquad u_\varphi = -\,c_n\,z_n(k_S r)\,\frac{dP_n}{d\theta}.
\end{equation}
Throughout, $dP_n/d\theta = P_n^1(\cos\theta)$ in the Condon--Shortley
convention, which follows from $P_n^1(u) = -\sqrt{1-u^2}\,P_n'(u)$ and
$dP_n/d\theta = -\sin\theta\,P_n'(u)$, the two coinciding because
$\sin\theta\ge0$ on $[0,\pi]$.  The combination $P_n^1(\cos\theta)/\sin\theta =
-P_n'(\cos\theta)$ is finite on the axis, with limit $-n(n+1)/2$.

\subsection{Traction on a plane of constant \texorpdfstring{$z$}{z}}\label{sec:traction}

The state used in \S\ref{sec:rt} is $(\bu,\bta_3)$ with
$\bta_3 = \bsig\cdot\zhat$.  From
$\bsig = \lambda(\nabla\cdot\bu)\bI + \mu(\nabla\bu + \nabla\bu\tp)$,
\begin{equation}\label{eq:tau3}
  (\bsig\cdot\zhat)_i = \lambda\,(\nabla\cdot\bu)\,\delta_{iz}
   + \mu\left(\partial_z u_i + \partial_i u_z\right),
\end{equation}
so only the $z$-row and $z$-column of the displacement gradient are required,
and $\nabla\cdot\bu$ is known in closed form from \eqref{eq:div}.

For $m=0$ the field is axisymmetric with no azimuthal component, so
$\sigma_{r\varphi} = \sigma_{\theta\varphi} = 0$ and the strain reduces to
\begin{equation}\label{eq:strain}
\begin{aligned}
  \varepsilon_{rr} &= \partial_r u_r, &
  \varepsilon_{\theta\theta} &= \frac{\partial_\theta u_\theta + u_r}{r}, \\
  \varepsilon_{\varphi\varphi} &= \frac{u_r + \cot\theta\,u_\theta}{r}, &
  \varepsilon_{r\theta} &= \tfrac12\!\left[\frac{\partial_\theta u_r}{r}
     + \partial_r u_\theta - \frac{u_\theta}{r}\right].
\end{aligned}
\end{equation}
Each term is finite on the axis: $u_\theta$ carries $dP_n/d\theta\sim\sin\theta$,
so $\cot\theta\,u_\theta$ is regular and is evaluated through
$P_n^1/\sin\theta$ rather than by dividing.  Legendre's equation removes the
second angular derivative,
\begin{equation}\label{eq:d2P}
  \frac{d^2P_n}{d\theta^2} = -\cot\theta\,\frac{dP_n}{d\theta} - n(n+1)P_n ,
\end{equation}
so no function beyond $P_n$ and $P_n^1$ is needed.  With
$\bsig = \lambda\,\mathrm{tr}(\varepsilon)\bI + 2\mu\varepsilon$ in the
orthonormal spherical basis, and $\rhat\cdot\zhat = \cos\theta$,
$\hat{\boldsymbol\theta}\cdot\zhat = -\sin\theta$,
$\hat{\boldsymbol\varphi}\cdot\zhat = 0$,
\begin{equation}\label{eq:tau3sph}
  \bta_3 = \left(\sigma_{rr}\cos\theta - \sigma_{r\theta}\sin\theta\right)\rhat
   + \left(\sigma_{r\theta}\cos\theta - \sigma_{\theta\theta}\sin\theta\right)
     \hat{\boldsymbol\theta}.
\end{equation}
The radial derivatives in \eqref{eq:strain} follow from
\eqref{eq:uL}--\eqref{eq:uM} with
$z_n'' = \big(n(n+1)/x^2 - 1\big)z_n - (2/x)z_n'$, the spherical Bessel
equation.

%=====================================================================
\section{The plane-wave representation}\label{sec:weyl}

\subsection{The identity}

Every outgoing multipole is a superposition of plane waves.  For $z > 0$,
\begin{equation}\label{eq:weyl}
  h^{(1)}_n(Kr)\,Y_n^m(\rhat)
   = \frac{1}{2\pi K\,i^{\,n}}
     \iint dk_x\,dk_y\;
     e^{\,i\bk\cdot\br}\;
     \frac{Y_n^m(\khat)}{k_z},
   \qquad \bk = (k_z,k_x,k_y),
\end{equation}
with $k_z$ on the branch \eqref{eq:branch}.  For $z < 0$ replace $k_z$ by
$-k_z$ in $\bk$ and in $\khat$, so that the exponential is
$e^{i(k_xx + k_yy + k_z|z|)}$ in either half-space.  At $n=0$, \eqref{eq:weyl}
is the classical Weyl identity for $h_0$.

The angular weight is the multipole's own harmonic \emph{evaluated at the
plane-wave direction}.  This is what makes the spectrum closed-form: no
transform of the field is involved.

\subsection{The harmonic must be the solid form}\label{sec:solid}

Outside the light circle $k_z$ is imaginary, so $\khat$ is a complex direction
and its polar angle is complex.  $Y_n^m$ must therefore be evaluated in a form
that is analytic in its arguments --- the solid harmonic, a homogeneous
degree-zero ratio of polynomials in $(k_x,k_y,k_z)$.  A representation through
real polar angles is defined only inside the light circle and silently discards
the evanescent band, which is most of the domain of integration.

In practice this means evaluating $P_n$, $P_n'$ and $P_n''$ at the complex
argument $k_z/K$, which \eqref{eq:angexplicit} requires; all three are
polynomials and take a complex argument without qualification.

%=====================================================================
\section{The three channel spectra}\label{sec:spectra}

\subsection{Polarisations}

Write $q = |\kperp|$, let $\psi_k$ be the azimuth of $\kperp$, and let
$k_z^{\pm}$ denote $+k_z$ below the sphere and $-k_z$ above it.  At each
$\kperp$ the scattered field is a sum of exactly three plane waves: one $P$ at
$k_{z,P}$ and two shear at $k_{z,S}$.  Their polarisations are
\begin{equation}\label{eq:pols}
  \bep_P = \frac{(k_z^{\pm},\,k_x,\,k_y)}{k_P},
  \qquad
  \bep_{SV} = \frac{(-q,\;k_z^{\pm}\cos\psi_k,\;k_z^{\pm}\sin\psi_k)}{k_S},
  \qquad
  \bep_{SH} = (0,\,-\sin\psi_k,\,\cos\psi_k),
\end{equation}
each of unit length, with $\bep_{SV}$ and $\bep_{SH}$ orthogonal to $\bk$.
They depend on the direction only, not on $n$ or $m$.

\subsection{The angular function and its surface gradient}

Let $A_{nm}$ be the angular function of the potential --- $P_n(\cos\theta)$ at
$m=0$, $P_n^1(\cos\theta)\cos\varphi$ at $m=1$ --- evaluated at the
\emph{plane-wave direction}, so that $\cos\theta \to k_z^{\pm}/K$,
$\sin\theta \to q/K$ and $\varphi \to \psi_k$.  The two transverse vector
harmonics built from it are
\begin{equation}\label{eq:vecharm}
  \nabla_{\!s}A = \frac{\partial A}{\partial\theta}\,\hat{\boldsymbol\theta}
    + \frac{1}{\sin\theta}\frac{\partial A}{\partial\varphi}\,\hat{\boldsymbol\varphi},
  \qquad
  \rhat\times\nabla_{\!s}A
    = -\frac{1}{\sin\theta}\frac{\partial A}{\partial\varphi}\,\hat{\boldsymbol\theta}
      + \frac{\partial A}{\partial\theta}\,\hat{\boldsymbol\varphi},
\end{equation}
with $\hat{\boldsymbol\theta}$ and $\hat{\boldsymbol\varphi}$ at the plane-wave
direction identified with $\bep_{SV}$ and $\bep_{SH}$ respectively.

\subsection{The spectra}

With the normalisations of \eqref{eq:weyl} cancelling, and writing
$D = 2\pi K\,i^{\,n} k_z$ with $K = k_P$ for the $L$-type and $k_S$ for the
other two,
\begin{equation}\label{eq:spectra}
\boxed{\;
\begin{aligned}
  \text{$L$-type } (a_n):\quad
    &\frac{i k_P}{D}\;A_{nm}\;\bep_P,\\[4pt]
  \text{$N$-type } (b_n):\quad
    &\frac{i k_S}{D}\left[\frac{\partial A_{nm}}{\partial\theta}\,\bep_{SV}
      + \frac{1}{\sin\theta}\frac{\partial A_{nm}}{\partial\varphi}\,\bep_{SH}\right],\\[4pt]
  \text{$M$-type } (c_n):\quad
    &\frac{-1}{D}\left[-\frac{1}{\sin\theta}\frac{\partial A_{nm}}{\partial\varphi}\,\bep_{SV}
      + \frac{\partial A_{nm}}{\partial\theta}\,\bep_{SH}\right],
\end{aligned}\;}
\end{equation}
each multiplied by its Mie coefficient and summed over $n$.  The three
constants $ik_P$, $ik_S$ and $-1$ are independent of $n$ and of $m$.

Explicitly, with $u = k_z^{\pm}/K$ and $s = q/K$,
\begin{equation}\label{eq:angexplicit}
\begin{aligned}
  m=0:\quad & A = P_n(u), &
    & \frac{\partial A}{\partial\theta} = -s\,P_n'(u), &
    & \frac{1}{\sin\theta}\frac{\partial A}{\partial\varphi} = 0,\\[4pt]
  m=1:\quad & A = -s\,P_n'(u)\cos\psi_k, &
    & \frac{\partial A}{\partial\theta} = \big[{-u}P_n'(u) + s^2P_n''(u)\big]\cos\psi_k, &
    & \frac{1}{\sin\theta}\frac{\partial A}{\partial\varphi} = P_n'(u)\sin\psi_k .
\end{aligned}
\end{equation}

\subsection{Three remarks}

\paragraph{The $L$-type is a multiplier; the others are not.}
$\bu = \nabla\phi$ gives the $P$ amplitude immediately.
$\nabla\times\nabla\times(\br\psi)$ carries an explicit position vector, and
acting on a single plane-wave component,
\begin{equation}\label{eq:obstruct}
  \nabla\times\nabla\times\!\left(\br\,e^{i\bk\cdot\br}\right)
   = e^{i\bk\cdot\br}\big[k_S^2\br - \bk(\bk\cdot\br)\big] - 2i\,\bk\,e^{i\bk\cdot\br},
\end{equation}
so the $\br$-dependence cancels only under the $\bk$ integration and not term
by term.  The results in \eqref{eq:spectra} are nonetheless exact.

\paragraph{Each shear family feeds \emph{both} transverse polarisations.}
That is the content of \eqref{eq:vecharm}, and it is why \eqref{eq:spectra}
cannot be written as one amplitude per family.  At $m=0$ the azimuthal
derivative vanishes, each family collapses onto a single polarisation, and the
structure is invisible --- which is the case most often written down.

\paragraph{The $M$-type constant carries no factor of $k$.}
$\nabla\times\nabla\times(\br\,\cdot)$ has one derivative more than
$\nabla\times(\br\,\cdot)$, so potentials normalised alike give amplitudes
differing by a power of $k_S$: the constants are $ik_S$ and $-1$ respectively.

%=====================================================================
\section{Evaluating the representation}\label{sec:eval}

The spectra \eqref{eq:spectra} are what $\bR$ and $\bT$ are built from, and
need no integration.  Reconstructing the field from them --- to verify the
construction, or to obtain the field at a point --- requires the integral in
\eqref{eq:weyl}, and two reductions make it straightforward.

\subsection{The azimuthal integral, in closed form}

For $m=0$ the only azimuthal dependence of the integrand is the plane-wave
phase and the polarisation.  With $\rho = \sqrt{x^2+y^2}$ and $\varphi$ its
azimuth,
\begin{equation}\label{eq:azim}
  \int_0^{2\pi}\!\! e^{iq\rho\cos(\psi_k-\varphi)}\,d\psi_k = 2\pi J_0(q\rho),
  \qquad
  \int_0^{2\pi}\!\!\cos\psi_k\,e^{iq\rho\cos(\psi_k-\varphi)}\,d\psi_k
    = 2\pi i\cos\varphi\,J_1(q\rho),
\end{equation}
and similarly with $\sin\psi_k \to \sin\varphi$.  The double integral collapses
to a one-dimensional integral in $q$.  The vertical component of the $P$ and
$SV$ fields pairs with $J_0$ and the horizontal with $J_1$; the $SH$ field,
having no vertical component, carries only $J_1$.

\subsection{The branch point, removed by substitution}

The integrand carries $1/k_z$, which is singular at $q = K$.  The singularity
is integrable but a quadrature driven across it converges only algebraically.
Splitting the path at the light circle and substituting
\begin{equation}\label{eq:sommer}
  q = K\sin\theta \;\;(0\le\theta\le\tfrac{\pi}{2}), \qquad
  q = K\cosh t \;\;(t\ge0),
\end{equation}
gives, respectively,
\begin{equation}
  \frac{q\,dq}{k_z} = K\sin\theta\,d\theta,
  \qquad
  \frac{q\,dq}{k_z} = -\,i\,K\cosh t\,dt,
\end{equation}
so the singular factor cancels identically and both pieces are analytic.  This
is the standard Sommerfeld path.  The upper limit $t_{\max}$ is the one
convergence parameter; the evanescent branch carries $e^{-K\sinh t\,|z|}$, so
convergence is geometric and $q_{\max}$ of a few times $K$ suffices unless the
plane is very close to the sphere.

%=====================================================================
\section{Reflection and transmission}\label{sec:rt}

\subsection{One-sidedness}

Above the sphere there are no sources, so every scattered wave crossing
$z = -z_p$ travels in $-z$: the field there is purely up-going.  Below
$z = +z_p$ it is purely down-going.  The representation \eqref{eq:weyl} with
$k_z \to -k_z$ above builds this in, so the spectra of \S\ref{sec:spectra} are
one-sided by construction and no separation step is required.

\subsection{The matrices}

Let the incident wave be of type $\iota\in\{P,SV,SH\}$ with unit amplitude, and
let $A^{(\iota)}_\nu(\kperp)$ be the amplitude of the outgoing wave of type
$\nu$ obtained from \eqref{eq:spectra} with the Mie coefficients for that
incidence, \S\ref{sec:incident}.  Then
\begin{equation}\label{eq:RT}
  \bR_{\nu\iota}(\kperp) = A^{(\iota)}_\nu(\kperp)\Big|_{z=-z_p},
  \qquad
  \bT_{\nu\iota}(\kperp) = \delta_{\nu\iota}\,(2\pi)^2\delta^2(\kperp-\kperp^0)
     + A^{(\iota)}_\nu(\kperp)\Big|_{z=+z_p} .
\end{equation}
The delta in $\bT$ is the unscattered incident wave \eqref{eq:incspec}: it
travels in $+z$, so it appears on the lower plane and not the upper one, which
is why $\bR$ is purely scattered.  In practice the delta is kept symbolically
and only the scattered part is computed; if a finite lateral grid is used it
becomes $\delta_{\kperp\kperp^0}/\Delta^2$ with $\Delta$ the grid spacing in
$\kperp$.

Both matrices are $3\times3$ in polarisation at each $\kperp$, and both carry
the $e^{ik_zz_p}$ phase of the plane they are referred to; changing $z_p$
multiplies them by a diagonal phase and nothing else, which is a useful check.

\subsection{State form}

If the response is wanted on the six-component state $(\bu,\bta_3)$ rather than
on mode amplitudes --- as it is when comparing against a depth-marching
scheme --- each mode contributes a column
\begin{equation}\label{eq:state}
  \begin{pmatrix}\bu\\[2pt]\bta_3\end{pmatrix}_\nu
  = \begin{pmatrix}
      \bep_\nu\\[4pt]
      i\big[\lambda(\bk\cdot\bep_\nu)\zhat + \mu\big(k_z\bep_\nu + (\bep_\nu\!\cdot\!\zhat)\bk\big)\big]
    \end{pmatrix},
\end{equation}
obtained by substituting a single plane wave into \eqref{eq:tau3}.  The six
columns --- three up, three down --- are independent at every $\kperp$, so a
state on a plane determines its up- and down-going content uniquely.  With
displacement alone they would not be: six three-vectors in three dimensions are
necessarily dependent, and the traction is what separates the two senses.

%=====================================================================
\section{Summary of the construction}

\begin{enumerate}
\item Solve the Mie boundary conditions at $r = a$ for $a_n$, $b_n$, $c_n$.
\item Fix $\kperp$; take $k_{z,P}$ and $k_{z,S}$ on the branch
      \eqref{eq:branch}, reversing their sign above the sphere.
\item Evaluate $P_n$, $P_n'$ and $P_n''$ at $k_z/k_P$ and $k_z/k_S$ ---
      complex arguments, polynomial form --- and assemble the angular function
      and its two surface derivatives, \eqref{eq:angexplicit}.
\item Form the polarisation-resolved amplitudes \eqref{eq:spectra}, weight by
      $a_n$, $b_n$, $c_n$, and sum over $n$.  Both shear families contribute to
      both transverse polarisations unless $m=0$.
\item Assemble $\bR$ and $\bT$ from \eqref{eq:RT}, or the state columns from
      \eqref{eq:state}.
\end{enumerate}

\noindent
No spatial grid, no transform, and no lateral truncation enters at any stage.
The only quadrature is the optional reconstruction of \S\ref{sec:eval}, whose
single parameter is the upper limit of a geometrically convergent integral.

\paragraph{On the azimuthal reduction.}  The closed forms \eqref{eq:azim} are
stated for $m=0$, where the only azimuthal dependence is the plane-wave phase.
At $m=1$ the angular function contributes its own $\cos\psi_k$, which multiplies
the polarisations' $\cos\psi_k$ and $\sin\psi_k$ and produces $J_0$, $J_1$ and
$J_2$ terms.  That bookkeeping is avoidable: the $\psi_k$ integrand is smooth
and $2\pi$-periodic, so a uniform rule converges spectrally and the same
quadrature serves every $m$.  Nothing in \eqref{eq:spectra} depends on which
route is taken --- the reduction is only for reconstructing the field, not for
forming the response.

\end{document}
